\chapter*{Kata Pengantar}
\addcontentsline{toc}{chapter}{Kata Pengantar}
\markboth{Kata Pengantar}{Kata Pengantar}

Edisi pertama buku ini tumbuh dari satu masalah sederhana: koleksi kitab ajar yang
saya kumpulkan berisi lebih dari sepuluh ribu halaman, tetapi halaman-halaman itu
tidak pernah terbaca berurutan. Membuka Rudin pada bab tiga tanpa pernah menyentuh
definisi supremum hanya menghasilkan frustrasi. Yang kurang bukan bahan, melainkan
urutan.

Sebelum naskah ini, ada versi tiga belas bab. Setelah dibaca ulang, dua kelemahan
kelihatan. Pertama, tiga belas bab terlalu cepat sehingga bagian tersulit hanya lewat
sekilas. Kedua, buku itu hanya bisa dibaca satu jenis pembaca: mahasiswa yang sudah
terbiasa dengan notasi formal. Siswa SMP yang penasaran tidak punya pintu masuk, dan
dosen tidak menemukan catatan mengajar di dalamnya.

Versi ini memperbaiki keduanya. Cakupannya diperluas menjadi enam bagian dan tiga puluh
dua bab, dengan materi yang di versi tiga belas bab tidak ada: himpunan dan
kardinalitas, deret tak hingga, sistem persamaan linear dan determinan, integral lipat
serta teorema lapangan vektor, persamaan diferensial dan sistem dinamik, regresi,
generalisasi dan regularisasi model. Setiap bab kini disusun berlapis.

Lapisan pertama adalah kotak \emph{Untuk pembaca muda}. Isinya analogi benda
sehari-hari, tanpa satu pun simbol matematika, ditutup satu pertanyaan yang membuat
penasaran. Lapisan kedua adalah badan bab itu sendiri: definisi, teorema, bukti,
contoh hitung, dan latihan, seperti buku ajar biasa. Lapisan ketiga adalah kotak
\emph{Catatan untuk dosen}: urutan mengajar yang saya sarankan, miskonsepsi yang
hampir selalu muncul, dan satu aktivitas kelas. Lapisan keempat adalah kotak
\emph{Untuk peneliti}: apa yang benar-benar terbukti, apa yang hanya dipopulerkan,
sumber primernya, dan pertanyaan yang masih terbuka.

Keempat lapisan tidak saling menggantikan: pembaca muda dapat membaca satu bab hanya
dari kotak oranye dan ringkasannya, mahasiswa membaca semuanya kecuali kotak peneliti,
dosen memakai kotak hijau sebagai rencana pertemuan, dan peneliti memulai dari kotak
abu-abu. Satu buku, empat jalur.

Lapisan itu mengatur \emph{kedalaman}. Yang mengatur \emph{cara masuk} adalah empat
kotak kedua: \emph{Bacakan keras} untuk pembaca yang paham dengan mendengar,
\emph{Kerjakan dengan tangan} untuk yang paham dengan bergerak dan meraba,
\emph{Resep langkah} untuk yang paham dengan menghafal prosedur lebih dahulu, dan
\emph{Tantangan} untuk yang paham justru ketika soalnya belum jelas. Keduanya hidup
bersama: satu bab yang sama, dengan enam pintu.

Tiga prinsip penulisan tidak berubah. \emph{Tidak ada klaim tanpa alat uji}: setiap
teorema disertai bukti atau pengakuan jujur bahwa bukti dilewati beserta lokasinya;
setiap angka berasal dari eksekusi nyata yang pembaca dapat ulangi. \emph{Notasi
diperkenalkan sebelum dipakai}. \emph{Kejujuran evaluasi adalah bagian dari
matematika}: bab-bab terakhir tidak berhenti pada cara membangun model, melainkan
pada cara mengetahui apakah klaim tentang model itu terbukti.

Satu hal yang sengaja tidak ada di buku ini: kunci jawaban lengkap. Memberi angka
akhir tanpa jejak pemeriksaan akan melanggar prinsip pertama. Sebagai gantinya,
Lampiran~C memberi kriteria penilaian dan alat verifikasi untuk setiap jenis latihan,
sehingga pembaca dapat memeriksa dirinya sendiri dengan cara yang sama seperti
penulis memeriksa naskah ini.

Buku ini bukan pengganti kitab sumbernya, melainkan pintu masuk yang tertata: jika
satu bab membuat pembaca membuka Rudin, Axler, Boyd, atau makalah Athalye lalu membaca
bagian yang ditunjuk, buku ini sudah berhasil. Kesalahan yang tersisa milik penulis,
bukan kitab-kitab itu.

\enlargethispage{3\baselineskip}
\vspace{0.35cm}
\begin{flushright}
Afiq Andico Pangimpian\\
Oktober 2026
\end{flushright}
